% sec_ub2.tex --- input into Section 6: "The upper bound in the regime r > d".
% Relies on the host for: \eqref{eq:triples}, Lemma~\ref{lem:norm},
% Lemma~\ref{lem:wlog}, Theorem~\ref{thm:main}, Sections~\ref{sec:family}
% and~\ref{sec:comp}, and the macros \Rad, \vtwo.
% The top-block offset is denoted E := d(a+1), avoiding a clash with the
% triple notation T(x,k) of \eqref{eq:triples}.

The proof is a chain of forcing steps. Two elementary \emph{pair
constraints}, coming from the multipliers $k = 1$ and $k = d+1$
(Lemma~\ref{lem:ub2-pairs}), compose into a $\pm 2d$ ladder showing that any
putative valid colouring is $2d$-periodic on an initial \emph{window}
(Lemma~\ref{lem:ub2-window}) and colour-\emph{flipped} on a \emph{top block}
of the same length (Lemma~\ref{lem:ub2-flip}); the window pattern
$\phi \colon \mathbb{Z}_{2d} \to \{0,1\}$ is then \emph{antipodal},
$\phi(\rho + d) = 1 - \phi(\rho)$ (Lemma~\ref{lem:ub2-antipodal}). Next,
three explicit families of conflict triples --- with multipliers $k = d+2$,
$k = d$, and a two-triple dichotomy using the consecutive multipliers
$k = i,\, i+1$ through a common cell --- forbid every \emph{$s$-plateau} of
$\phi$ (Lemma~\ref{lem:ub2-plateau}), and a purely combinatorial case
analysis upgrades this to \emph{$s$-antiperiodicity},
$\phi(m + s) = 1 - \phi(m)$ (Lemma~\ref{lem:ub2-rigidity}). Finally, a
$2$-adic computation shows that antipodality and $s$-antiperiodicity are
incompatible exactly when $\vtwo(s) \ne \vtwo(d)$ --- and this is precisely
the existence condition in this regime (Lemma~\ref{lem:2adic-b}) --- so no
valid colouring exists. The final obstruction is thus \emph{exactly} the
$2$-adic existence condition: the $s$-antiperiodic patterns are the unique
family of colourings escaping the plateau conflicts, and the existence
condition is precisely what rules them out.

\paragraph{Setup and standing assumption.}
Fix $d \ge 2$ and $a > d$ in the regime $r > d$: by Lemma~\ref{lem:norm},
under the existence condition $\vtwo(a) \le \vtwo(a-d)$ we have
$r \in [d+1, 2d-1]$, so we write
\[
a \;=\; 2dq + d + s, \qquad s := r - d \in [1, d-1], \qquad q \ge 0,
\]
and set $c = a - d$. Put
\[
N \;:=\; (a+3)d + a - (d-1) \;=\; (d+1)a + 2d + 1,
\qquad
E \;:=\; d(a+1),
\qquad
W \;:=\; a + d + 1 .
\]
Then $N = E + W$: indeed
$(a+3)d + a - (d-1) = ad + 3d + a - d + 1 = (d+1)a + 2d + 1$ and
$E + W = d(a+1) + a + d + 1 = (d+1)a + 2d + 1 = N$. Moreover
$a \equiv d + s \pmod{2d}$, and $a \ge d + 1$ (as $q \ge 0$ and $s \ge 1$),
so $W = a + d + 1 \ge 2d + 2 > 2d$. (For $q \ge 1$ one has $a \ge 3d + 1$;
the proofs below use only $a \ge d + 1$.) By \eqref{eq:triples}, the
solutions of $x_1 + a x_2 - x_3 = c$ in $[1,N]^3$ are exactly the triples
$T(x,k) = (x,\, k,\, x + a(k-1) + d)$ with $x, k \ge 1$ and
$x + a(k-1) + d \le N$ (entries may repeat).

Suppose, for contradiction, that $\chi \colon [1,N] \to \{0,1\}$ is a valid
colouring, i.e.\ admits no monochromatic solution; by
Lemma~\ref{lem:wlog} we may and do assume $\chi(1) = 0$. This
\emph{standing assumption} --- a valid colouring $\chi$ of $[1,N]$ with
$\chi(1) = 0$, together with the quantities $N$, $E$, $W$ above --- remains
in force throughout Lemmas~\ref{lem:ub2-pairs}--\ref{lem:ub2-plateau};
Lemmas~\ref{lem:ub2-rigidity} and~\ref{lem:2adic-b} are self-contained
statements about patterns on $\mathbb{Z}_{2d}$ and $2$-adic valuations.

\begin{lemma}[pair constraints]\label{lem:ub2-pairs}
\begin{enumerate}
\item[(a)] $\chi(d+1) = 1$.
\item[(b)] (R0) There is no $x \in [1, N-d]$ with $\chi(x) = \chi(x+d) = 0$.
\item[(c)] (R1) There is no $x \in [1, W]$ with $\chi(x) = \chi(x+E) = 1$.
\end{enumerate}
\end{lemma}

\begin{proof}
(b) For $x \in [1, N-d]$, $T(x,1) = (x, 1, x+d)$ is a solution
($x_3 = x + a \cdot 0 + d$). If $\chi(x) = \chi(x+d) = 0$, then, since
$\chi(1) = 0$, the triple $T(x,1)$ is monochromatic of colour $0$ ---
contradiction.

(a) Apply (b) with $x = 1$: $T(1,1) = (1, 1, 1+d)$ is a solution
($1 + d \le N$), and $\chi(1) = 0$, so $\chi(d+1) = 0$ would make it
monochromatic; hence $\chi(d+1) = 1$.

(c) For $x \in [1, W]$, $T(x, d+1) = (x, d+1, x+E)$ is a solution:
$x + a\bigl((d+1)-1\bigr) + d = x + ad + d = x + d(a+1) = x + E$, and
$x + E \le W + E = N$. If $\chi(x) = \chi(x+E) = 1$, then, since
$\chi(d+1) = 1$ by (a), the triple $T(x, d+1)$ is monochromatic of colour
$1$ --- contradiction.
\end{proof}

We refer to the constraints of Lemma~\ref{lem:ub2-pairs}(b) and (c) as R0
and R1.

\begin{lemma}[window periodicity]\label{lem:ub2-window}
There is a function $\phi \colon \mathbb{Z}_{2d} \to \{0,1\}$ such that
$\chi(x) = \phi(x \bmod 2d)$ for every $x \in [1, W]$. Moreover
$\phi(1) = 0$ and $\phi(d+1) = 1$.
\end{lemma}

\begin{proof}
\emph{Claim U (up-step): if $\chi(x) = 0$ and $1 \le x \le W - 2d$, then
$\chi(x + 2d) = 0$.} Indeed: (i) $\chi(x) = 0$ and $x + d \le W \le N$
imply $\chi(x + d) = 1$ by R0. (ii) $x + d \in [1, W]$, so R1 gives
$\chi(x + d + E) = 0$. (iii) Apply R0 at position $x + d + E$: it lies in
$[1, N-d]$ because $(x + d + E) + d = x + 2d + E \le W + E = N$ (using
$x \le W - 2d$); since $\chi(x + d + E) = 0$, R0 gives
$\chi(x + 2d + E) = 1$. (iv) Apply R1 at position $x + 2d \in [1, W]$:
since $\chi(x + 2d + E) = 1$, R1 forces $\chi(x + 2d) = 0$. This proves
Claim U.

\emph{Claim D (down-step): if $\chi(y) = 0$ and $2d + 1 \le y \le W$, then
$\chi(y - 2d) = 0$.} Indeed: (i) R0 at $y - d$ (note $y - d \ge d + 1 \ge 1$
and $(y - d) + d = y \le W \le N$): $\chi(y) = 0$ forces $\chi(y - d) = 1$.
(ii) R1 at $y - d \in [1, W]$: $\chi(y - d + E) = 0$. (iii) R0 at
$y - 2d + E$ (note $y - 2d + E \ge E + 1 \ge 1$ and
$(y - 2d + E) + d = y - d + E \le W + E = N$): $\chi(y - d + E) = 0$ forces
$\chi(y - 2d + E) = 1$. (iv) R1 at $y - 2d \in [1, W]$:
$\chi(y - 2d + E) = 1$ forces $\chi(y - 2d) = 0$. This proves Claim D.

\emph{Class constancy.} Fix $\rho \in [1, 2d]$. The members of the residue
class $\{x \equiv \rho \pmod{2d}\}$ inside $[1, W]$ are
$z_j = \rho + 2dj$, $j = 0, \dots, \ell$, where
$\ell := \lfloor (W - \rho)/(2d) \rfloor$; the list is nonempty since
$\rho \le 2d < W$ (recall $W \ge 2d + 2$). Suppose some member $z_i$ has
$\chi(z_i) = 0$. Applying Claim D repeatedly from $z_i$ (each intermediate
member is $\le W$, and $\ge 2d + 1$ while $j \ge 1$) gives $\chi(z_j) = 0$
for all $j \le i$; in particular $\chi(z_0) = \chi(\rho) = 0$. Applying
Claim U repeatedly from $z_0$ (the step from $z_j$ to $z_{j+1}$ needs
$z_j \le W - 2d$, which holds exactly when $z_{j+1} \le W$, i.e.\ whenever
$z_{j+1}$ is a member) gives $\chi(z_j) = 0$ for all $j \le \ell$. Hence
either all members are $0$ or none is: $\chi$ is constant on each residue
class inside $[1, W]$. Define $\phi(\rho)$ to be that constant. Then
$\phi(1) = \chi(1) = 0$, and $\phi(d+1) = \chi(d+1) = 1$ by
Lemma~\ref{lem:ub2-pairs}(a) (both positions lie in $[1, W]$).
\end{proof}

\begin{lemma}[flip on the top block]\label{lem:ub2-flip}
For every $x \in [1, W]$: $\chi(E + x) = 1 - \chi(x)$. Consequently
$\chi(E + x) = 1 - \phi(x \bmod 2d)$ on the top block
$[E+1, E+W] = [E+1, N]$, which together with
Lemma~\ref{lem:ub2-window} determines $\chi$ on $[1, W] \cup [E+1, N]$.
\end{lemma}

\begin{proof}
\emph{Case $\chi(x) = 1$}: R1 at $x \in [1, W]$ directly forbids
$\chi(x + E) = 1$, so $\chi(E + x) = 0$.

\emph{Case $\chi(x) = 0$ and $x \le a + 1$}: (i) $x + d \le a + d + 1 = W
\le N$, so R0 gives $\chi(x + d) = 1$. (ii) R1 at $x + d \in [1, W]$ gives
$\chi(x + d + E) = 0$. (iii) Apply R0 at position $x + E$: it lies in
$[1, N-d]$ because $(x + E) + d = x + d + E \le (a+1) + d + E = N$; since
$\chi(x + d + E) = 0$, R0 forbids $\chi(x + E) = 0$; hence
$\chi(x + E) = 1$.

\emph{Case $\chi(x) = 0$ and $a + 2 \le x \le W$}: (i)
$x - d \ge a + 2 - d \ge 1$ (as $a \ge d + 1$), and R0 at $x - d$ forbids
$\chi(x - d) = 0 = \chi(x)$; hence $\chi(x - d) = 1$. (ii) R1 at
$x - d \in [1, W]$ gives $\chi(x - d + E) = 0$. (iii) R0 at $x - d + E$
(note $(x - d + E) + d = x + E \le W + E = N$) forbids $\chi(x + E) = 0$;
hence $\chi(x + E) = 1$.

In all cases $\chi(E + x) = 1 - \chi(x)$; by Lemma~\ref{lem:ub2-window},
$\chi(x) = \phi(x \bmod 2d)$ for $x \in [1, W]$.
\end{proof}

\begin{lemma}[antipodality of the window pattern]\label{lem:ub2-antipodal}
$\phi(\rho + d) = 1 - \phi(\rho)$ for every $\rho \in \mathbb{Z}_{2d}$.
\end{lemma}

\begin{proof}
It suffices to prove the identity for $\rho \in [1, d]$, since the pairs
$\{\rho, \rho + d\}$ partition $\mathbb{Z}_{2d}$. For such $\rho$, both
positions $\rho$ and $\rho + d$ lie in $[1, 2d]$, and $2d \le W$, so
$\chi(\rho) = \phi(\rho)$ and $\chi(\rho + d) = \phi(\rho + d)$.

\emph{Not both $0$}: if $\phi(\rho) = \phi(\rho + d) = 0$, then
$\chi(\rho) = \chi(\rho + d) = 0$ with $\rho \le N - d$, contradicting R0.

\emph{Not both $1$}: if $\phi(\rho) = \phi(\rho + d) = 1$, then by
Lemma~\ref{lem:ub2-flip}, $\chi(E + \rho) = 0$ and
$\chi(E + \rho + d) = 0$ (both $\rho$ and $\rho + d$ lie in $[1, W]$). The
position $E + \rho$ lies in $[1, N - d]$, since
$E + \rho + d \le E + 2d \le E + W = N$. So R0 is violated at $E + \rho$
--- contradiction.

Hence exactly one of $\phi(\rho)$, $\phi(\rho + d)$ is $0$, i.e.\
$\phi(\rho + d) = 1 - \phi(\rho)$.
\end{proof}

\begin{lemma}[plateau-conflict triples]\label{lem:ub2-plateau}
Under the standing assumption (a valid colouring $\chi$ of $[1,N]$ with
$\chi(1) = 0$) and with the objects of
Lemmas~\ref{lem:ub2-window}--\ref{lem:ub2-antipodal}, call $m \in [1, d+1]$
an \emph{$s$-plateau} of colour $\kappa$ if
$\phi(m) = \phi(m + s) = \kappa$ (indices modulo $2d$; note
$m + s \le 2d$, so no wrap occurs). Then:
\begin{enumerate}
\item[(V)] there is no $s$-plateau of colour $1 - \phi(2)$;
\item[(IV)] there is no $s$-plateau of colour $1 - \phi(d)$;
\item[(Cb)] there is no pair $(i, m)$ with $i \in [1, d-1]$,
$\phi(i) \ne \phi(i+1)$, and $m$ an $s$-plateau of colour $\phi(i+1)$.
\end{enumerate}
\end{lemma}

\begin{proof}
Throughout the proof we use: $\chi(x) = \phi(x \bmod 2d)$ on $[1, W]$
(Lemma~\ref{lem:ub2-window}); $\chi(E + x) = 1 - \phi(x \bmod 2d)$ for
$x \in [1, W]$ (Lemma~\ref{lem:ub2-flip});
$\phi(\rho + d) = 1 - \phi(\rho)$ (Lemma~\ref{lem:ub2-antipodal}); and
$a \equiv d + s \pmod{2d}$, whence
$\phi(a + m) = \phi(d + s + m) = 1 - \phi(m + s)$.

(V): Suppose $\phi(m) = \phi(m + s) = 1 - \phi(2) =: \kappa$ for some
$m \in [1, d+1]$. Consider the triple $T(m, d+2) = (m,\, d+2,\, E + a + m)$.
It is a solution: $m + a\bigl((d+2) - 1\bigr) + d = m + (d+1)a + d
= m + E + a$, and $E + a + m \le E + a + d + 1 = N$; all entries are
$\ge 1$. Colours: $\chi(m) = \phi(m) = \kappa$ (as $m \le d + 1 \le W$);
$\chi(d+2) = \phi(d+2) = 1 - \phi(2) = \kappa$ (as $d + 2 \le W$);
$\chi\bigl(E + (a + m)\bigr) = 1 - \phi(a + m) = 1 - \bigl(1 - \phi(m + s)\bigr)
= \phi(m + s) = \kappa$ (as $a + m \le a + d + 1 = W$, so
Lemma~\ref{lem:ub2-flip} applies). The triple is monochromatic ---
contradiction.

(IV): Suppose $\phi(m) = \phi(m + s) = 1 - \phi(d)$ for some
$m \in [1, d+1]$. Consider the triple $T(a + m, d) = (a + m,\, d,\, E + m)$.
It is a solution: $(a + m) + a(d - 1) + d = da + d + m = E + m \le N$; all
entries are $\ge 1$. Colours:
$\chi(a + m) = \phi(a + m) = 1 - \phi(m + s) = \phi(d)$ (as $a + m \le W$);
$\chi(d) = \phi(d)$ (as $d \le W$);
$\chi(E + m) = 1 - \phi(m) = 1 - \bigl(1 - \phi(d)\bigr) = \phi(d)$ (as
$m \le W$). The triple is monochromatic of colour $\phi(d)$ ---
contradiction.

(Cb): Suppose $i \in [1, d-1]$ with $\phi(i) \ne \phi(i+1) =: \kappa$, and
$m \in [1, d+1]$ with $\phi(m) = \phi(m + s) = \kappa$. Let
$y := (d + 1 - i)a + m$. Then
$1 \le 2a + 1 \le y \le da + d + 1 = E + 1 \le N$. Consider two triples
through $y$:
\begin{align*}
\alpha &:= T(y,\, i+1) = (y,\; i+1,\; E + a + m),
&&\text{a solution, since } y + ai + d = (d+1)a + m + d = E + a + m \le N
\ (m \le d+1);\\
\beta &:= T(y,\, i) = (y,\; i,\; E + m),
&&\text{a solution, since } y + a(i-1) + d = da + m + d = E + m \le N .
\end{align*}
Colours of the known cells: $\chi(i) = \phi(i) = 1 - \kappa$ and
$\chi(i+1) = \phi(i+1) = \kappa$ (as $i + 1 \le d \le W$);
$\chi(E + a + m) = 1 - \phi(a + m) = \phi(m + s) = \kappa$ (as
$a + m \le W$); $\chi(E + m) = 1 - \phi(m) = 1 - \kappa$ (as $m \le W$).
Now $\chi(y)$ is either $\kappa$ --- making $\alpha$ monochromatic of
colour $\kappa$ --- or $1 - \kappa$ --- making $\beta$ monochromatic of
colour $1 - \kappa$. Either way a monochromatic solution exists ---
contradiction.
\end{proof}

\begin{lemma}[rigidity: $\phi$ must be $s$-antiperiodic]\label{lem:ub2-rigidity}
Let $\phi \colon \mathbb{Z}_{2d} \to \{0,1\}$ be antipodal
($\phi(\rho + d) = 1 - \phi(\rho)$) with $\phi(1) = 0$, and suppose that
none of the configurations (V), (IV), (Cb) of
Lemma~\ref{lem:ub2-plateau} occurs. Then $\phi(m + s) = 1 - \phi(m)$ for
\emph{all} $m \in \mathbb{Z}_{2d}$; that is, $\phi$ is
\emph{$s$-antiperiodic}.
\end{lemma}

\begin{proof}
First, if $\phi(m^*) = \phi(m^* + s)$ for some $m^* \in \mathbb{Z}_{2d}$
(take the representative $m^* \in [1, 2d]$), then there is an $s$-plateau
with index in $[1, d+1]$: if $m^* \in [1, d+1]$, take $m = m^*$; if
$m^* \in [d+2, 2d]$, take $m = m^* - d \in [2, d]$ --- by antipodality
$\phi(m) = 1 - \phi(m^*)$ and $\phi(m + s) = 1 - \phi(m^* + s)$, so
$\phi(m) = \phi(m + s)$: an in-range $s$-plateau (of the opposite colour).

Now suppose, for contradiction, that some $s$-plateau exists; by the
previous paragraph the set $K$ of colours of in-range $s$-plateaus
($m \in [1, d+1]$) is nonempty. Since (V) does not occur, no in-range
plateau has colour $1 - \phi(2)$; hence $K = \{\kappa\}$ with
$\kappa := \phi(2)$. Since (IV) does not occur, $1 - \phi(d) \notin K$, so
$\phi(d) = \kappa$.

\emph{Case $\kappa = 1$}: then $\phi(1) = 0 \ne 1 = \phi(2)$, so $i = 1$
(legal, since $d \ge 2$ gives $1 \le d - 1$) is a transition with
$\phi(i + 1) = \kappa$, and an in-range plateau of colour $\kappa$ exists
--- configuration (Cb) occurs, a contradiction.

\emph{Case $\kappa = 0$}: then $\phi(1) = \phi(2) = \phi(d) = 0$.
\emph{Sub-claim: $\phi \equiv 0$ on $[1, d]$.} For $d \in \{2, 3\}$ this is
immediate --- the values $\phi(1) = \phi(2) = \phi(d) = 0$ already exhaust
$[1, d]$ --- so the branch that follows is vacuous there. Otherwise, let
$j$ be minimal in $[1, d]$ with $\phi(j) = 1$; then $3 \le j \le d - 1$.
Since $\phi(d) = 0$, let $i$ be minimal in $[j, d-1]$ with
$\phi(i + 1) = 0$; by minimality $\phi(i) = 1$ (for $i = j$ directly; for
$i > j$, all of $\phi(j), \dots, \phi(i)$ equal $1$). Then
$i \in [1, d-1]$, $\phi(i) \ne \phi(i+1) = 0 = \kappa$, and an in-range
plateau of colour $\kappa$ exists --- configuration (Cb) occurs, a
contradiction. This proves the sub-claim: $\phi \equiv 0$ on $[1, d]$, and
hence, by antipodality, $\phi \equiv 1$ on $[d+1, 2d]$. But then
$m = d + 1 \in [1, d+1]$ satisfies $\phi(d + 1) = 1$ and
$\phi(d + 1 + s) = 1$ (since $d + 2 \le d + 1 + s \le 2d$, because
$1 \le s \le d - 1$): an in-range plateau of colour $1$, so
$1 \in K = \{0\}$ --- a contradiction.

Both cases are impossible, so no $s$-plateau exists anywhere in
$\mathbb{Z}_{2d}$, i.e.\ $\phi(m + s) = 1 - \phi(m)$ for all $m$.
\end{proof}

\begin{lemma}[$2$-adic obstruction and existence criterion]\label{lem:2adic-b}
Let $d \ge 2$, $q \ge 0$, and $a = 2dq + d + s$.
\begin{enumerate}
\item[(a)] For $s \in [1, d-1]$, the existence condition
$\vtwo(a) \le \vtwo(a - d)$ is equivalent to $\vtwo(s) \ne \vtwo(d)$.
Moreover, the same computation covers $s = d$ (i.e.\ $r = 2d$): there the
existence condition always fails, re-proving Lemma~\ref{lem:norm}(i), so
that under the existence condition the regime $r \in [d+1, 2d]$ reduces to
$s \in [1, d-1]$.
\item[(b)] If $\vtwo(s) \ne \vtwo(d)$, then no
$\phi \colon \mathbb{Z}_{2d} \to \{0,1\}$ can be simultaneously antipodal
and $s$-antiperiodic.
\end{enumerate}
\end{lemma}

\begin{proof}
(a) Write $v := \vtwo(d)$ and $w := \vtwo(s)$, so that $a = 2dq + (d + s)$
and $a - d = 2dq + s$, and note $\vtwo(2dq) \ge v + 1$ whenever $q \ge 1$
(if $q = 0$ the term is absent).

If $w < v$: $\vtwo(d + s) = w$ and $\vtwo(2dq) \ge v + 1 > w$, so
$\vtwo(a) = w$; likewise $\vtwo(a - d) = \vtwo(2dq + s) = w$. Hence
$\vtwo(a) = \vtwo(a - d)$: the existence condition holds.

If $w > v$: $\vtwo(d + s) = v < \vtwo(2dq)$, so $\vtwo(a) = v$; and
$\vtwo(a - d) \ge \min(v + 1,\, w) \ge v + 1 > \vtwo(a)$: the existence
condition holds (strictly).

If $w = v$: write $d = g d'$, $s = g s'$ with $g = \gcd(d, s)$; then
$\vtwo(g) = v$ and $d'$, $s'$ are both odd, so $\vtwo(d + s) \ge v + 1$
and hence $\vtwo(a) \ge v + 1$; but
$\vtwo(a - d) = \vtwo(2dq + s) = v$ (since
$\vtwo(2dq) \ge v + 1 > v = \vtwo(s)$). So $\vtwo(a) > \vtwo(a - d)$: the
existence condition fails. This proves the equivalence for
$s \in [1, d-1]$.

Finally, the computation in the case $w = v$ applies verbatim to $s = d$:
there $g = d$ and $d' = s' = 1$ are both odd, so again
$\vtwo(a) > \vtwo(a - d)$. Thus $r = 2d$ never satisfies the existence
condition, and, under the existence condition, the regime
$r \in [d+1, 2d]$ reduces to $s \in [1, d-1]$.

(b) Suppose $\phi$ is antipodal and $s$-antiperiodic:
$\phi(x + d) = 1 - \phi(x)$ and $\phi(x + s) = 1 - \phi(x)$ for all
$x \in \mathbb{Z}_{2d}$. Composing the two,
$\phi\bigl(x + (d + s)\bigr) = \phi(x)$ and
$\phi\bigl(x + (d - s)\bigr) = \phi(x)$: both $d + s$ and $d - s$ are
periods of $\phi$. The set $P$ of periods of $\phi$ is a subgroup of
$\mathbb{Z}_{2d}$, so $P$ contains the subgroup generated by $d - s$ and
$d + s$, which is generated by $h := \gcd(d - s,\, d + s,\, 2d)$. With
$g = \gcd(d, s)$ and $d = g d'$, $s = g s'$:
$\gcd(d - s,\, d + s) = g \gcd(d' - s',\, d' + s')$, and
$\gcd(d' - s',\, d' + s')$ divides both $2d'$ and $2s'$, hence divides
$2 \gcd(d', s') = 2$. Since $\vtwo(s) \ne \vtwo(d)$, the integers $d'$ and
$s'$ have opposite parities, so $d' + s'$ is odd, so
$\gcd(d' - s',\, d' + s')$ is odd, hence equals $1$. Thus
$\gcd(d - s,\, d + s) = g$, and $h = \gcd(g, 2d) = g$ (as $g \mid d$). So
$g \in P$, and since $g \mid d$, also $d \in P$: $\phi(x + d) = \phi(x)$
for all $x$. But antipodality says $\phi(x + d) = 1 - \phi(x)$ ---
contradiction.
\end{proof}

\begin{theorem}[upper bound in the regime $r > d$]\label{thm:ub2}
For every $d \ge 2$, $q \ge 0$ and $s \in [1, d-1]$ with
$\vtwo(s) \ne \vtwo(d)$, put $a := 2dq + d + s$, so that $r = d + s > d$ ---
equivalently, consider any $a$ for which $\Rad(a;\, a-d)$ exists
and $r > d$. Then every $2$-colouring of $[1, N]$,
$N = (a+3)d + a - (d-1)$, contains a monochromatic solution of
$x_1 + a x_2 - x_3 = a - d$; that is,
\[
\Rad(a;\, a - d) \;\le\; (a+3)d + a - (d-1).
\]
\end{theorem}

\begin{proof}
Suppose $\chi \colon [1, N] \to \{0,1\}$ has no monochromatic solution;
by Lemma~\ref{lem:wlog} we may assume $\chi(1) = 0$, and the standing
assumption of this section is in force.
Lemmas~\ref{lem:ub2-pairs}--\ref{lem:ub2-antipodal} produce the antipodal
window pattern $\phi$ with $\phi(1) = 0$. Since $\chi$ is assumed valid,
none of the configurations (V), (IV), (Cb) of
Lemma~\ref{lem:ub2-plateau} can occur (each produces an explicit
monochromatic solution inside $[1, N]$). By
Lemma~\ref{lem:ub2-rigidity}, $\phi$ is $s$-antiperiodic. By
Lemma~\ref{lem:2adic-b}(a) and (b), this contradicts the existence
condition: $\vtwo(s) \ne \vtwo(d)$ holds, yet no antipodal
$s$-antiperiodic pattern on $\mathbb{Z}_{2d}$ exists. Hence no such
$\chi$ exists, and $\Rad(a;\, a - d) \le N$.
\end{proof}

Since $r > d$ gives $\min(r-1,\, d-1) = d - 1$, Theorem~\ref{thm:ub2} is
exactly the upper bound required by Theorem~\ref{thm:main} in this regime;
combined with the extremal family colouring of $[1, N-1]$ of
Section~\ref{sec:family}, it pins
$\Rad(a;\, a-d) = (a+3)d + a - \min(r-1,\, d-1)$ in the whole regime
$r > d$.

\begin{remark}
All linear range conditions used in the proofs of
Lemmas~\ref{lem:ub2-pairs}--\ref{lem:ub2-plateau} (each cited triple
satisfies $x_3 = x_1 + a(x_2 - 1) + d$ with all entries in $[1, N]$) were
verified symbolically above and, in addition, mechanically replayed at the
boundary values of all free parameters for all $1769$ in-scope pairs
$(a, d)$ with $d \le 16$, $a \le 400$ (including $q = 0$). Likewise, the
case analysis of Lemma~\ref{lem:ub2-rigidity} was corroborated by
exhaustive enumeration: for every $d \le 16$ and every $s$, among all
$2^{d-1}$ antipodal patterns $\phi$ with $\phi(1) = 0$, those avoiding the
three configurations of Lemma~\ref{lem:ub2-plateau} are exactly the
$s$-antiperiodic ones --- and, in accordance with
Lemma~\ref{lem:2adic-b}, none exists when $\vtwo(s) \ne \vtwo(d)$. See
Section~\ref{sec:comp}.
\end{remark}
